\documentclass[10pt, oneside]{article}
\usepackage{amsmath, amsthm, amssymb, calrsfs, wasysym, verbatim, bbm, color, graphics, geometry}
\usepackage{graphicx}
\graphicspath{ {./assets/} }
\geometry{tmargin=.75in, bmargin=.75in, lmargin=.75in, rmargin = .75in}

\newcommand{\R}{\mathbb{R}}
\newcommand{\C}{\mathbb{C}}
\newcommand{\Z}{\mathbb{Z}}
\newcommand{\N}{\mathbb{N}}
\newcommand{\Q}{\mathbb{Q}}
\newcommand{\Cdot}{\boldsymbol{\cdot}}

\newtheorem{thm}{Theorem}
\newtheorem{defn}{Definition}
\newtheorem{conv}{Convention}
\newtheorem{rem}{Remark}
\newtheorem{lem}{Lemma}
\newtheorem{cor}{Corollary}
\usepackage[backend=biber,style=authoryear]{biblatex}
\addbibresource{reference.bib}

\title{Summary Report: Candidate literary paper to be used}
\author{Alejandro Ouslan}
\date{2026-02-08}

\begin{document}

\maketitle
\tableofcontents

\vspace{.25in}

\section{Paper 1: Optimum government size and economic growth in case of Indian states: Evidence from panel threshold model}

\subsection{Summary}
This paper investigates the relationship between optimum government size and economic growth using data of
Indian states during 1990-91 to 2017–18. Our results derived from panel threshold regression model show a
positive and significant impact of government size on economic growth within the estimated thresholds for both
aggregate and sub-panels based on income and regions. Once the government size moves above the upper
threshold level, then its impact declines and turns to be insignificant. Thus, our findings suggest the policymakers
for maintaining the government size within the thresholds limit. \cite{akram2020optimum}

\subsection{Data sources}
\begin{itemize}
	\item Indian states during 1990-91 to 2017–18
	\item Gross State Domestic Products (GSDP, 2011–12) and expenditures (such total expenditure, social sector expenditure
	\item revenue expenditure, and capital expenditure) data collected from Economic and Political Weekly Research Foundation (EPWRF).
\end{itemize}

\subsection{Methodology}
theoretical model which is proposed by Barro (1990, 1991).
$$
	Y(t) = AK(t)^\beta G(t)^{1-\beta}, \quad 0< \alpha < 1
$$

The thresholds model is as follows:
$$
	\delta y_{i,t} = \alpha_) + \pi_i t + \beta_i y_{t-1} + \sum^k_{j=1} \psi_{i,j} \delta y_{i,t-1} + \epsilon_{i,t}
$$

\subsection{Conclusions}
\begin{itemize}
	\item \textbf{Objective:} The study addresses the literature gap regarding sub-national economic growth by examining the relationship between optimum government size and growth across major Indian states (1990--91 to 2017--18).
	\item \textbf{Methodology:} A panel threshold regression model was employed, identifying two significant thresholds for the aggregate panel and most sub-panels.
	\item \textbf{Optimal Thresholds (Aggregate):} Government size has a positive and significant effect on economic growth specifically when it lies between the lower and upper thresholds of $4.85\%$ and $5.09\%$.
	\item \textbf{Category Variations:}
	      \begin{itemize}
		      \item \textbf{NSC States:} Estimated optimal threshold range is $4.55\%$--$6.07\%$.
		      \item \textbf{SC States:} Estimated optimal threshold range is significantly higher at $34.70\%$--$38.55\%$.
	      \end{itemize}
	\item \textbf{Diminishing Returns:} Once government size exceeds the upper threshold, the coefficient of growth declines and becomes statistically insignificant.
	\item \textbf{Expenditure Composition:} Findings remain consistent when disaggregated into revenue, capital, and social sector expenditures, showing positive growth associations within threshold limits.
	\item \textbf{Policy Recommendations:} State governments should utilize fiscal rules and budgetary discipline to maintain an ideal size, ensuring that Keynesian deficit spending remains within the "optimum" range to avoid unfavorable growth consequences.
\end{itemize}

\section{Paper 2: Does too much government spending depress the
  economic development of transition economies?
  Evidences from dynamic panel threshold analysis}

\subsection{Summary}
This paper investigates the relationship between government size and economic growth and
determines the optimal level of government spending to maximize economic growth. The paper
applies a dynamic panel data analysis based upon a threshold model to test the threshold effect
of government spending in 26 transition economies over the period spanning 1993–2016.
According to the analysis results, government expenditures have a threshold effect on economic
growth, and there is a non-linear relationship depicted as an Armey curve in these transition
economies. The findings indicate that a government size above the threshold government
spending level adversely affects economic growth, while a government size below the threshold
level has a positive effect. Furthermore, there is a statistically significant relationship between the
two variables above and below that optimal level, even if we divide the sample into developed
and developing countries. Our findings suggest that governments in transition economies should
consider optimal government size at around the estimated threshold level to support sustainable
economic growth.\cite{aydin2019does}
\subsection{Methodology}
h$$
	Y_{it} = \alpha_0 + \alpha_i initial_{it} + \alpha_2 Gov_{it} + \beta x_{it} + \epsilon_{it}
$$
Where:
\begin{itemize}
	\item $Y_{it}$: Real GDP per capita growth rate in country $i$ at time $t$.
	\item $\alpha_0$: Constant term.
	\item $initial_{it}$: Initial level of income in country $i$ at time $t$.
	\item $\alpha_1$: Coefficient measuring the effect of initial income on economic growth.
	\item $Gov_{it}$: Government spending level in country $i$ at time $t$.
	\item $\alpha_2$: Coefficient capturing the impact of government spending on economic growth.
	\item $x_{it}$: Vector of other macroeconomic control variables affecting economic growth.
	\item $\beta$: Vector of coefficients associated with the control variables.
	\item $\varepsilon_{it}$: White noise error term.
\end{itemize}
\subsection{Conclusions}
\begin{itemize}
	\item The relationship between government expenditures and economic growth is widely debated, but has rarely been studied in the context of transition economies moving from centrally planned to free market systems.
	\item This study empirically examines threshold effects of government size on economic growth in 25 transition economies from 1993 to 2016.
	\item Using a threshold autoregressive (TAR) approach, the study avoids assuming a specific functional form and identifies optimal government size more accurately.
	\item Results show a non-linear relationship: government expenditures positively affect growth below a threshold, but this effect weakens and becomes negative (though insignificant) beyond it.
	\item The estimated threshold levels are 11.67 \% of GDP for developing economies and 17.54 \% for developed economies.
	\item Government intervention supports growth through education, health, inequality reduction, and institutional development, but excessive intervention leads to diminishing returns.
	\item Beyond the optimal size, government expansion may hinder growth due to taxation burdens, inefficiencies, corruption, rent-seeking, and crowding out of private investment.
	\item Transition economies require some degree of government intervention (e.g., liberalization, privatization, legal reforms) to support market functionality during the transition process.
	\item Sustainable growth depends on effective and limited government intervention; exceeding the optimal size does more harm than good.
	\item The study also highlights that government impact extends beyond expenditure size to include regulatory, legal, and bureaucratic interventions.
\end{itemize}
\printbibliography
\end{document}
